\chapter{विटिग और C=C बनाने वाली अन्य अभिक्रियाएँ}\label{ch:b2:wittig}

मादा घरेलू मक्खी 23 कार्बनों वाले एक हाइड्रोकार्बन से नरों को आकर्षित करती है, जिसमें केवल एक द्वि-आबंध है, कार्बन
9 और 10 के बीच, (\textit{Z}) ज्यामिति के साथ। उसे बड़ी मात्रा में बनाने के लिए रसायनज्ञ को वह एक
द्वि-आबंध ठीक वहीं, और उसी ज्यामिति के साथ, रखना होता है: कोई विलोपन उसे कई
स्थितियों पर बिखेर देता और दोनों ज्यामितियाँ देता। यह अध्याय द्वि-आबंध दूसरे ढंग से बनाता है, किसी कार्बोनिल समूह पर दो
टुकड़ों को जोड़कर, जिससे द्वि-आबंध वहीं प्रकट होता है जहाँ \ce{C=O} था।
% ledger: use:muscalure

\begin{recall}
\Cref{ch:b2:enolates-aldol}: \omterm{def:b2:enolates-aldol:aldol}{ऐल्डोल संघनन}। \Cref{ch:b2:alkene-redox}: विषाक्त उत्प्रेरक पर किसी
ऐल्काइन का (\textit{Z})-ऐल्कीन तक हाइड्रोजनन। प्रथम वर्ष का खंड: E1, E2 और ज़ैत्सेव का
नियम, द्वि-अणुक प्रतिस्थापन, ग्रीन्यार अभिकर्मक और ध्रुवता प्रतिलोमन,
(\textit{Z})/(\textit{E}) वर्णनक। विद्यालय का खंड: परमाणु मितव्ययिता।
\end{recall}

\section{फ़ॉस्फ़ोनियम लवण और इलाइड}

\begin{definition}[फ़ॉस्फ़ोनियम इलाइड]\label{def:b2:wittig:ylide}
\emph{फ़ॉस्फ़ोनियम लवण}\index{फ़ॉस्फ़ोनियम लवण} $\mathrm{R_3P^+{-}CH_2R'\ X^-}$ किसी फ़ॉस्फ़ीन, प्रायः
ट्राइफ़ेनिलफ़ॉस्फ़ीन \ce{Ph3P}, द्वारा किसी हैलोऐल्केन के द्वि-अणुक प्रतिस्थापन से बनता है। फ़ॉस्फ़ोरस से बँधे कार्बन से
एक हाइड्रोजन हटाने पर \emph{फ़ॉस्फ़ोनियम इलाइड}\index{फ़ॉस्फ़ोनियम इलाइड} मिलता है,
जो सन्निकट विपरीत आवेशों वाला एक उदासीन अणु है, $\mathrm{Ph_3P^+{-}CH^-{-}R'}$, जिसे $\mathrm{Ph_3P{=}CHR'}$ भी लिखा जाता है।
\emph{स्थायीकृत इलाइड}\index{स्थायीकृत इलाइड} उस कार्बन पर कोई इलेक्ट्रॉन-आकर्षी समूह
(एस्टर, कीटोन, \omterm{def:b2:amines:nitrile}{नाइट्राइल}) धारण करता है जो ऋणात्मक आवेश को विस्थानीकृत करता है।
\end{definition}

\begin{omfigure}
\setchemfig{atom sep=1.9em}
\schemestart
\chemfig{Ph_3P} \+ \chemfig{CH_3-I}
\arrow{->[S$_\mathrm{N}$2]}[,0.9]
\chemfig{Ph_3P^{+}-CH_3}\;\chemfig{I^{-}}
\arrow{->[\ce{BuLi}][$-$ \ce{BuH}, \ce{LiI}]}[,1.3]
\chemfig{Ph_3P^{+}-\chemabove{C}{\scriptstyle -}H_2}
\arrow{<->}[,0.8]
\chemfig{Ph_3P=CH_2}
\schemestop
\par\vspace{6pt}

{\small मेथिलिडीनट्राइफ़ेनिलफ़ॉस्फ़ोरेन, सबसे सरल इलाइड। ट्राइफ़ेनिलफ़ॉस्फ़ीन आयडोमेथेन से आयोडाइड को
विस्थापित करता है; ब्यूटिललिथियम मेथिल समूह से एक प्रोटॉन हटाता है। दाईं ओर की दोनों संरचनाएँ
एक ही अणु को लिखने के दो ढंग हैं: आवेश-पृथक्कृत संरचना वास्तविकता के अधिक निकट है।}
\end{omfigure}

\begin{proposition}[इलाइड बनाना]\label{prop:b2:wittig:ylide-acidity}
\omterm{def:b2:wittig:ylide}{फ़ॉस्फ़ोनियम लवण} फ़ॉस्फ़ोरस के बगल वाले कार्बन पर किसी ऐल्केन से कहीं अधिक अम्लीय है। सरल ऐल्किल
\omterm{def:b2:wittig:ylide}{फ़ॉस्फ़ोनियम लवण} को फिर भी शुष्क, अक्रिय परिस्थितियों में अत्यंत प्रबल क्षारक (ब्यूटिललिथियम, सोडियम हाइड्राइड, सोडियम ऐमाइड)
चाहिए; ऐसे लवण का, जिसके \ce{CH2} पर एक कार्बोनिल समूह भी हो, किसी दुर्बल क्षारक
(सोडियम कार्बोनेट, सोडियम हाइड्रॉक्साइड) से विप्रोटॉनन हो जाता है, और बने \omterm{def:b2:wittig:ylide}{स्थायीकृत इलाइड} को पृथक और संगृहीत किया जा सकता है।
\end{proposition}

\begin{proof}[तर्क]
कार्बऋणायन के बगल वाला धनात्मक फ़ॉस्फ़ोरस अपने आवेश से और बड़े फ़ॉस्फ़ोरस परमाणु की
ध्रुवणीयता से उसे स्थायी करता है: सरल इलाइड एक स्थायीकृत कार्बऋणायन है, परंतु \omterm{def:b2:enolates-aldol:enolate}{ईनोलेट} से कम,
इसीलिए प्रबल क्षारक चाहिए। सन्निकट कार्बोनिल ऋणात्मक आवेश को ऑक्सीजन पर फैला देता है, जैसे \omterm{def:b2:enolates-aldol:enolate}{ईनोलेट} में
(\cref{prop:b2:enolates-aldol:alpha-acidity}), फ़ॉस्फ़ोरस के अतिरिक्त: संयुग्मी अम्ल इतना
अम्लीय हो जाता है कि दुर्बल क्षारक पर्याप्त हो, और इलाइड, जो कहीं कम क्षारकीय है, अब जल या वायु से अभिक्रिया नहीं करता।
\end{proof}

\section{विटिग अभिक्रिया}

\begin{definition}[विटिग अभिक्रिया]\label{def:b2:wittig:wittig}
किसी ऐल्डिहाइड या कीटोन के साथ किसी \omterm{def:b2:wittig:ylide}{फ़ॉस्फ़ोनियम इलाइड} की \emph{विटिग अभिक्रिया}\index{विटिग अभिक्रिया}
एक ऐल्कीन और ट्राइफ़ेनिलफ़ॉस्फ़ीन ऑक्साइड देती है:
\begin{equation*}
\mathrm{Ph_3P{=}CHR + R'CHO \longrightarrow R'CH{=}CHR + Ph_3P{=}O}.
\end{equation*}
यह एक \emph{ऑक्साफ़ॉस्फ़ेटेन}\index{ऑक्साफ़ॉस्फ़ेटेन} से होकर जाती है, जो फ़ॉस्फ़ोरस, दो
कार्बनों और ऑक्सीजन का चतुःसदस्यी वलय है, और तब बनता है जब इलाइड कार्बन कार्बोनिल कार्बन से आबंधित होता है जबकि कार्बोनिल ऑक्सीजन
फ़ॉस्फ़ोरस से आबंधित होता है।
\end{definition}

\begin{omfigure}
\setchemfig{atom sep=1.9em}
\schemestart
\chemfig{Ph_3P=CH-R} \+ \chemfig{O=CH-R'}
\arrow{->[योग]}[,1.0]
\chemfig{Ph_3P?-[6]CH(-[4]R)-[0]CH(-[0]R')-[2]O?}
\schemestop

\medskip
\schemestart
\arrow{->[विलोपन]}[,1.1]
\chemfig{R-CH=CH-R'} \+ \chemfig{Ph_3P=O}
\schemestop
\par\vspace{6pt}

{\small \omterm{def:b2:wittig:wittig}{विटिग अभिक्रिया}। इलाइड कार्बन और कार्बोनिल कार्बन आपस में आबंधित होते हैं जबकि ऑक्सीजन
फ़ॉस्फ़ोरस से आबंधित होता है, एक ही चरण में: \omterm{def:b2:wittig:wittig}{ऑक्साफ़ॉस्फ़ेटेन}। फिर इसका वलय दूसरी ओर से खुलता है, और
\ce{P-C} तथा \ce{C-O} आबंध टूटते हैं: \ce{C=C} आबंध दोनों भूतपूर्व अभिक्रियाशील कार्बनों के बीच बनता है, और
ट्राइफ़ेनिलफ़ॉस्फ़ीन ऑक्साइड का \ce{P=O} आबंध उसी समय बनता है।}
\end{omfigure}

\begin{proposition}[ठीक स्थान पर द्वि-आबंध]\label{prop:b2:wittig:regiospecific}
\omterm{def:b2:wittig:wittig}{विटिग अभिक्रिया} में नया \ce{C=C} आबंध भूतपूर्व कार्बोनिल कार्बन को भूतपूर्व इलाइड कार्बन से जोड़ता है,
और कहीं नहीं।
\end{proposition}

\begin{proof}[तर्क]
बनने और टूटने वाले आबंध केवल चतुःसदस्यी वलय के हैं: इलाइड कार्बन से कार्बोनिल
कार्बन, फिर \ce{C-P} और \ce{C-O} आबंध। कोई हाइड्रोजन नहीं खिसकता और शृंखला के अनुदिश कोई कार्बधनायन या कार्बऋणायन
नहीं बनता, अतः द्वि-आबंध खिसक नहीं सकता। इसके विपरीत कोई विलोपन कई पड़ोसी कार्बनों के
हाइड्रोजनों के बीच चयन करता है (ज़ैत्सेव का नियम), और कोई E1 अभिक्रिया पुनर्विन्यस्त हो सकती है।
\end{proof}

\begin{proposition}[प्रेरक बल]\label{prop:b2:wittig:driving}
ट्राइफ़ेनिलफ़ॉस्फ़ीन ऑक्साइड के \ce{P=O} आबंध का बनना \omterm{def:b2:wittig:wittig}{विटिग अभिक्रिया} को अत्यंत
अनुकूल और अनुत्क्रमणीय बनाता है।
\end{proposition}

\begin{proof}[तर्क]
अभिक्रिया एक \ce{C=O} और एक \ce{P-C} आबंध (इलाइड में \ce{P=C} मुख्यतः एक
\ce{P+-C-} एकल आबंध है) के बदले एक \ce{C=C} और एक \ce{P=O} आबंध देती है। किसी फ़ॉस्फ़ीन ऑक्साइड का \ce{P=O} आबंध
अत्यंत प्रबल है, खोए गए \ce{P-C} आबंध से कहीं प्रबल, और खोए गए \ce{C=O} की आंशिक क्षतिपूर्ति
बने \ce{C=C} से होती है: संतुलन स्पष्ट रूप से ढलान पर है। फ़ॉस्फ़ोरस अभिक्रिया का
ऑक्सीजन-अवशोषक है।
\end{proof}

\section{त्रिविमवरणात्मकता}

\begin{proposition}[ऐल्कीन की ज्यामिति]\label{prop:b2:wittig:stereo}
किसी ऐल्डिहाइड के साथ अस्थायीकृत इलाइड (\ce{Ph3P=CHR}, R एक ऐल्किल समूह) मुख्यतः (\textit{Z})-ऐल्कीन देता है,
लवण-मुक्त परिस्थितियों में। \omterm{def:b2:wittig:ylide}{स्थायीकृत इलाइड} (\ce{Ph3P=CH-COOR}, \ce{Ph3P=CH-COR}) मुख्यतः
(\textit{E})-ऐल्कीन देता है। केवल ऐरिल समूह से \omterm{def:b2:wittig:ylide}{स्थायीकृत इलाइड} मिश्रण देता है।
\end{proposition}

\admitted

ये प्रेक्षण हैं। इनकी सामान्य व्याख्या \omterm{def:b2:wittig:wittig}{विटिग अभिक्रिया} को दरों और स्थायित्वों के बीच
प्रतिस्पर्धा के रूप में पढ़ती है। अभिक्रियाशील, अस्थायीकृत इलाइड के साथ \omterm{def:b2:wittig:wittig}{ऑक्साफ़ॉस्फ़ेटेन} तेज़ी से और
अनुत्क्रमणीय रूप से बनता है, कम भीड़ वाली संक्रमण अवस्था से होकर, जिसमें दोनों प्रतिस्थापी वलय की
एक ही ओर पहुँचते हैं; फिर यह उस व्यवस्था को बनाए रखते हुए टूटता है: गतिक नियंत्रण
(\textit{Z})-ऐल्कीन देता है। \omterm{def:b2:wittig:ylide}{स्थायीकृत इलाइड} के साथ पहला चरण धीमा और उत्क्रमणीय है;
\omterm{def:b2:wittig:wittig}{ऑक्साफ़ॉस्फ़ेटेन} आरंभिक पदार्थों से होकर परस्पर बदलते हैं, और वह जिसके प्रतिस्थापी
विपरीत ओर हैं, कम तनावग्रस्त, जीतता है: (\textit{E})-ऐल्कीन, ऊष्मागतिक नियंत्रण में।

\begin{omfigure}
\setchemfig{atom sep=1.9em}
\schemestart
\chemfig{Ph_3P=CH-CH_2CH_3} \+ \chemfig{O=CH-Ph}
\arrow{->}[,0.8]
\chemfig{Ph-[:-30]=[:30]-[:90]CH_2CH_3}
\schemestop

\medskip
\schemestart
\chemfig{Ph_3P=CH-COOEt} \+ \chemfig{O=CH-Ph}
\arrow{->}[,0.8]
\chemfig{Ph-[:-30]=[:30]-[:-30]COOEt}
\schemestop
\par\vspace{6pt}

{\small दो इलाइड, एक ऐल्डिहाइड। ऊपर: अस्थायीकृत प्रोपिलिडीन इलाइड मुख्यतः
(\textit{Z})-1-फ़ेनिलब्यूट-1-ईन देता है। नीचे: \omterm{def:b2:wittig:ylide}{स्थायीकृत इलाइड} मुख्यतः एथिल
(\textit{E})-3-फ़ेनिलप्रोपीनोएट (एथिल सिनेमेट) देता है। दोनों में बनने वाला ट्राइफ़ेनिलफ़ॉस्फ़ीन ऑक्साइड
नहीं दिखाया गया।}
\end{omfigure}

\section{हॉर्नर--वड्सवर्थ--एमन्स अभिक्रिया}

\begin{definition}[फ़ॉस्फ़ोनेट और HWE अभिक्रिया]\label{def:b2:wittig:hwe}
\emph{फ़ॉस्फ़ोनेट}\index{फ़ॉस्फ़ोनेट} किसी फ़ॉस्फ़ोनिक अम्ल का एस्टर है, $\mathrm{(RO)_2P({=}O){-}R'}$, जिसमें
एक \ce{P-C} आबंध हो। यह किसी ट्राइऐल्किल फ़ॉस्फ़ाइट की किसी हैलोऐल्केन के साथ माइकेलिस--आर्बुज़ोव अभिक्रिया से बनता है,
$\mathrm{P(OEt)_3 + R'Br \to (EtO)_2P(O)R' + EtBr}$। \emph{हॉर्नर--वड्सवर्थ--एमन्स अभिक्रिया}\index{हॉर्नर--वड्सवर्थ--एमन्स अभिक्रिया} (HWE) में,
अपने \ce{CH2} पर इलेक्ट्रॉन-आकर्षी समूह धारण करने वाले किसी फ़ॉस्फ़ोनेट का ऋणायन किसी ऐल्डिहाइड से जुड़ता है और एक
ऐल्कीन तथा एक डाइऐल्किल फ़ॉस्फ़ेट ऋणायन देता है।
\end{definition}

\begin{omfigure}
\setchemfig{atom sep=1.8em}
\schemestart
\chemfig{EtO-P(=[2]O)(-[6]OEt)-CH_2-COOEt}
\arrow{->[\ce{NaH}][$-$ \ce{H2}]}[,0.9]
\chemfig{EtO-P(=[2]O)(-[6]OEt)-\chemabove{C}{\scriptstyle -}H-COOEt}
\schemestop

\medskip
\schemestart
\arrow{->[\ce{PhCHO}]}[,0.9]
\chemfig{*6(=-=(-[:30]=[:-30]-[:30](=[:90]O)-[:-30]O-[:30]Et)-=-)}
\+
\chemfig{EtO-P(=[2]O)(-[6]O^{-})-OEt}
\schemestop
\par\vspace{6pt}

{\small बेन्ज़ैल्डिहाइड के साथ ट्राइएथिल फ़ॉस्फ़ोनोऐसीटेट की \omterm{def:b2:wittig:hwe}{हॉर्नर--वड्सवर्थ--एमन्स अभिक्रिया}। सोडियम
हाइड्राइड \omterm{def:b2:wittig:hwe}{फ़ॉस्फ़ोनेट} और एस्टर के बीच का हाइड्रोजन हटाता है; ऋणायन ऐल्डिहाइड से जुड़ता है,
और चतुःसदस्यी मध्यवर्ती \omterm{def:b2:wittig:wittig}{विटिग अभिक्रिया} की तरह टूटता है। उत्पाद एथिल
(\textit{E})-सिनेमेट है; उपोत्पाद डाइएथिल फ़ॉस्फ़ेट ऋणायन है, जो जल में विलेय है।}
\end{omfigure}

\begin{proposition}[रसायनज्ञ HWE अभिक्रिया को क्यों पसंद करते हैं]\label{prop:b2:wittig:hwe}
HWE अभिक्रिया मुख्यतः (\textit{E})-ऐल्कीन देती है, और उसका फ़ॉस्फ़ोरस उपोत्पाद, एक जल-विलेय
लवण, साधारण जलीय धुलाई से हट जाता है; \omterm{def:b2:wittig:hwe}{फ़ॉस्फ़ोनेट} ऋणायन संगत \omterm{def:b2:wittig:ylide}{स्थायीकृत इलाइड} से अधिक
नाभिकरागी भी है और कीटोनों से भी अभिक्रिया करता है।
\end{proposition}

\begin{proof}[तर्क]
\omterm{def:b2:wittig:hwe}{फ़ॉस्फ़ोनेट} ऋणायन एक स्थायीकृत कार्बऋणायन है, अतः उसका योग उत्क्रमणीय है और अभिक्रिया
ऊष्मागतिक नियंत्रण में चलती है: (\textit{E}), जैसे \omterm{def:b2:wittig:ylide}{स्थायीकृत इलाइडों} के साथ। \omterm{def:b2:wittig:ylide}{फ़ॉस्फ़ोनियम इलाइड} के विपरीत, यह
पूर्ण ऋणात्मक आवेश वहन करता है, जिसकी क्षतिपूर्ति किसी पड़ोसी धनायन से नहीं होती: यह अधिक अभिक्रियाशील है।
उपोत्पाद, \ce{(EtO)2PO2-}, एक आयनिक लवण है, जबकि ट्राइफ़ेनिलफ़ॉस्फ़ीन ऑक्साइड एक उदासीन कार्बनिक
ठोस है जिसे उत्पाद से वर्णलेखन या क्रिस्टलन द्वारा अलग करना पड़ता है।
\end{proof}

\section{विधि का चयन}

\begin{method}[विटिग विच्छेदन]\label{met:b2:wittig:wittig-disconnection}
\begin{enumerate}
  \item लक्ष्य \ce{C=C} आबंध को काटिए; एक कार्बन कार्बोनिल कार्बन बन जाता है, दूसरा इलाइड कार्बन
        (किसी हैलाइड से बना एक \omterm{def:b2:wittig:ylide}{फ़ॉस्फ़ोनियम लवण})।
  \item ऐसा करने के दो तरीकों में से उसे वरीयता दीजिए जिसका हैलाइड प्राथमिक हो (\ce{Ph3P} द्वारा
        द्वि-अणुक प्रतिस्थापन) और जिसका कार्बोनिल यौगिक ऐल्डिहाइड हो।
  \item वांछित ज्यामिति के अनुसार अभिकर्मक चुनिए: (\textit{Z}) के लिए अस्थायीकृत इलाइड,
        (\textit{E}) के लिए \omterm{def:b2:wittig:ylide}{स्थायीकृत इलाइड} या \omterm{def:b2:wittig:hwe}{फ़ॉस्फ़ोनेट} (HWE)।
  \item जाँचिए कि अणु में और कुछ भी प्रयुक्त क्षारक से अभिक्रिया न करे।
\end{enumerate}
\end{method}

\begin{method}[C=C बनाने वाली अभिक्रिया का चयन]\label{met:b2:wittig:choose-cc}
\begin{center}
\small
\begin{tabular}{@{}lll@{}}
\hline
विधि & \ce{C=C} की स्थिति & ज्यामिति \\
\hline
E2 या E1 विलोपन & ज़ैत्सेव का नियम, मिश्रण & अधिकांशतः (\textit{E}), मिश्रण \\
\omterm{def:b2:enolates-aldol:aldol}{ऐल्डोल संघनन} & \ce{C=O} से संयुग्मित & (\textit{E}) \\
विटिग, अस्थायीकृत इलाइड & भूतपूर्व \ce{C=O} पर & (\textit{Z}) \\
विटिग, \omterm{def:b2:wittig:ylide}{स्थायीकृत इलाइड}; HWE & भूतपूर्व \ce{C=O} पर & (\textit{E}) \\
ऐल्काइन, फिर विषाक्त उत्प्रेरक & भूतपूर्व त्रि-आबंध पर & (\textit{Z}) \\
\hline
\end{tabular}
\end{center}
\medskip
किसी \ce{C=C} पर समान आकार के दो टुकड़ों को जोड़ने के लिए \omterm{def:b2:wittig:wittig}{विटिग अभिक्रिया} और उसकी संबंधी अभिक्रियाएँ प्रायः
पहला विकल्प होती हैं; ओलेफ़िन मेटाथीसिस, जो दो ऐल्कीनों के सिरों का आदान-प्रदान करती है, तृतीय वर्ष के
खंड में आती है।
\end{method}

\begin{history}[वह इलाइड जो अभिकर्मक बन गया]
1950 के दशक के आरंभ में फ़ॉस्फ़ोरस इलाइडों का अध्ययन करते हुए गेओर्ग विटिग ने पाया कि मेथिलिडीनट्राइफ़ेनिलफ़ॉस्फ़ोरेन
बेन्ज़ोफ़ीनोन को 1,1-डाइफ़ेनिलएथीन और ट्राइफ़ेनिलफ़ॉस्फ़ीन ऑक्साइड में बदल देता है; अपने छात्र उलरिख
श्योलकोप्फ़ के साथ उन्होंने 1954 में इस प्रेक्षण को एक सामान्य विधि में बदल दिया। उद्योग ने इसे शीघ्र अपनाया,
अन्य उपयोगों के साथ विटामिन A के लिए, और विटिग ने 1979 का रसायन का नोबेल पुरस्कार साझा किया।
\end{history}

\begin{safety}{\ghs[5mm]{GHS02}\,\ghs[5mm]{GHS05}\,\ghs[5mm]{GHS07}\,\ghs[5mm]{GHS08}}
ब्यूटिललिथियम विलयन ज्वलनशील और संक्षारक हैं और जल से प्रचंड अभिक्रिया करते हैं; सोडियम हाइड्राइड
जल के साथ हाइड्रोजन मुक्त करता है। ट्राइफ़ेनिलफ़ॉस्फ़ीन हानिकारक और आँखों के लिए संक्षारक है, ट्राइएथिल
फ़ॉस्फ़ाइट ज्वलनशील और हानिकारक; ट्राइफ़ेनिलफ़ॉस्फ़ीन ऑक्साइड हानिकारक है।
% ledger: ghs:BuLi, ghs:PPh3, ghs:triethyl-phosphite, ghs:Ph3PO
\end{safety}

\section{अभ्यास}

\begin{exercise}[$\star$]\label{exo:b2:wittig:1}
ट्राइफ़ेनिलफ़ॉस्फ़ीन और एक हैलोऐल्केन से इलाइड \ce{Ph3P=CHCH3} का निर्माण लिखिए।
\end{exercise}

\begin{exercise}[$\star$]\label{exo:b2:wittig:2}
आयडोमेथेन से बने इलाइड के साथ बेन्ज़ैल्डिहाइड की \omterm{def:b2:wittig:wittig}{विटिग अभिक्रिया} के उत्पाद दीजिए।
\end{exercise}

\begin{exercise}[$\star$]\label{exo:b2:wittig:3}
स्थायीकृत, अस्थायीकृत या ऐरिल-स्थायीकृत के रूप में वर्गीकृत कीजिए: \ce{Ph3P=CHCOOEt}, \ce{Ph3P=CHCH3},
\ce{Ph3P=CHCOCH3}, \ce{Ph3P=CHC6H5}। किसे ब्यूटिललिथियम चाहिए?
\end{exercise}

\begin{exercise}[$\star$]\label{exo:b2:wittig:4}
ट्राइएथिल फ़ॉस्फ़ोनोऐसीटेट, सोडियम हाइड्राइड और प्रोपेनल का उत्पाद, उसकी ज्यामिति सहित, दीजिए।
\end{exercise}

\begin{exercise}[$\star\star$]\label{exo:b2:wittig:5}
मेथिलीनसाइक्लोहेक्सेन चाहिए। साइक्लोहेक्सेनोन से विटिग मार्ग की तुलना 1-मेथिलसाइक्लोहेक्सेन-1-ऑल के अम्ल-उत्प्रेरित
निर्जलीकरण से कीजिए।
\end{exercise}

\begin{exercise}[$\star\star$]\label{exo:b2:wittig:6}
\omterm{def:b2:wittig:wittig}{विटिग अभिक्रिया} द्वारा (\textit{E})-1,2-डाइफ़ेनिलएथीन (स्टिलबीन) का विच्छेदन कीजिए। शुद्ध (\textit{E})-स्टिलबीन के लिए विटिग मार्ग
ख़राब विकल्प क्यों है, और इसके बजाय आप क्या प्रयुक्त करेंगे?
\end{exercise}

\begin{exercise}[$\star\star$]\label{exo:b2:wittig:7}
प्रोपेनल एक ओर \ce{Ph3P=CHCH3} से और दूसरी ओर \ce{Ph3P=CHCOOEt} से अभिक्रिया करता है। उत्पाद और
उनकी मुख्य ज्यामिति दीजिए।
\end{exercise}

\begin{exercise}[$\star\star$]\label{exo:b2:wittig:8}
\ce{Ph3P=CH2} द्वारा साइक्लोहेक्सेनोन के मेथिलीनन की परमाणु मितव्ययिता परिकलित कीजिए।
\end{exercise}

\begin{exercise}[$\star\star$]\label{exo:b2:wittig:9}
एथिल ब्रोमोएथेनोएट के साथ ट्राइएथिल फ़ॉस्फ़ाइट की माइकेलिस--आर्बुज़ोव अभिक्रिया दो चरणों में लिखिए:
फ़ॉस्फ़ोरस द्वारा प्रतिस्थापन, फिर ब्रोमाइड द्वारा विऐल्किलन।
\end{exercise}

\begin{exercise}[$\star\star\star$]\label{exo:b2:wittig:10}
दो विटिग अभिक्रियाओं द्वारा (\textit{E})-ब्यूट-2-ईनडाइऐल
\ce{OHC-CH=CH-CHO} से 1,6-डाइफ़ेनिलहेक्सा-1,3,5-ट्राइईन का एक संश्लेषण प्रस्तावित कीजिए, और बताइए कि कौन-सा इलाइड और कितना।
\end{exercise}

\begin{exercise}[$\star\star\star$]\label{exo:b2:wittig:11}
(\textit{E})-4-फ़ेनिलब्यूट-3-ईन-2-ओन \omterm{def:b2:enolates-aldol:aldol}{ऐल्डोल संघनन} से या \omterm{def:b2:wittig:wittig}{विटिग अभिक्रिया} से बनाया जा सकता है।
दोनों मार्ग लिखिए और उनकी तुलना कीजिए।
\end{exercise}

\begin{exercise}[$\star\star\star$]\label{exo:b2:wittig:12}
(\textit{Z})-ऑक्ट-4-ईन तक दो मार्ग प्रस्तावित कीजिए, एक \omterm{def:b2:wittig:wittig}{विटिग अभिक्रिया} से और एक ऐल्काइन से होकर,
और बताइए कि हर एक कौन-से उपोत्पाद छोड़ता है।
\end{exercise}

\section{समस्या: विटिग द्वारा एक फ़ेरोमोन}

\begin{problem}[{सप्ताहांत समस्या --- किसी (Z)-ऐल्कीन का पश्च-संश्लेषण, अस्थायीकृत इलाइड और
उसकी वरणात्मकता, ट्राइफ़ेनिलफ़ॉस्फ़ीन ऑक्साइड की क़ीमत, और ऐल्काइन विकल्प}]
\label{pb:b2:wittig:1}
लक्ष्य: (\textit{Z})-ट्राइकोस-9-ईन, \ce{CH3(CH2)7CH=CH(CH2)12CH3}, आरंभ में वर्णित घरेलू मक्खी का
फ़ेरोमोन। मोलर द्रव्यमान (\unit{g/mol}): C 12.011, H 1.008, O 15.999, P 30.974, Br 79.904।
% ledger: use:muscalure, aw:C, aw:H, aw:O, aw:P, aw:Br, ghs:nonanal, ghs:BuLi

\medskip
\noindent\textbf{भाग I --- पश्च-संश्लेषण।}
\begin{enumerate}
  \item लक्ष्य का आण्विक सूत्र लिखिए और जाँचिए कि यह एक ऐल्कीन है।
  \item विटिग विच्छेदन कौन-सा आबंध काटता है?
  \item दो संभावित युग्म (कार्बोनिल यौगिक, \omterm{def:b2:wittig:ylide}{फ़ॉस्फ़ोनियम लवण}) दीजिए।
  \item नोनेनल + इलाइड युग्म के लिए कौन-सा हैलोऐल्केन \omterm{def:b2:wittig:ylide}{फ़ॉस्फ़ोनियम लवण} बनाता है?
  \item इस चरण के लिए प्राथमिक हैलोऐल्केन उपयुक्त क्यों है?
  \item \omterm{def:b2:wittig:ylide}{फ़ॉस्फ़ोनियम लवण} का बनना लिखिए।
\end{enumerate}

\medskip
\noindent\textbf{भाग II --- इलाइड और ज्यामिति।}
\begin{enumerate}[resume]
  \item कौन-सा क्षारक इस लवण का विप्रोटॉनन करता है? सोडियम हाइड्रॉक्साइड काम क्यों नहीं करेगा?
  \item क्या इलाइड स्थायीकृत है? इससे किस ज्यामिति का पूर्वानुमान होता है?
  \item नोनेनल के साथ बनने वाला \omterm{def:b2:wittig:wittig}{ऑक्साफ़ॉस्फ़ेटेन} लिखिए।
  \item अभिक्रिया शुष्क अवस्था में और अक्रिय गैस के अंतर्गत क्यों चलानी चाहिए?
  \item द्वि-आबंध ठीक कार्बन 9 और 10 के बीच ही क्यों पहुँचता है?
  \item इसके बजाय ट्राइकोसेन-9-ऑल का अम्ल-उत्प्रेरित निर्जलीकरण क्या देता?
\end{enumerate}

\medskip
\noindent\textbf{भाग III --- ट्राइफ़ेनिलफ़ॉस्फ़ीन ऑक्साइड।}
\begin{enumerate}[resume]
  \item इलाइड और नोनेनल से समग्र समीकरण लिखिए।
  \item लक्ष्य, नोनेनल, \omterm{def:b2:wittig:ylide}{फ़ॉस्फ़ोनियम लवण} और
        ट्राइफ़ेनिलफ़ॉस्फ़ीन ऑक्साइड के मोलर द्रव्यमान परिकलित कीजिए।
  \item विटिग चरण (इलाइड + ऐल्डिहाइड) की परमाणु मितव्ययिता परिकलित कीजिए।
  \item ट्राइफ़ेनिलफ़ॉस्फ़ीन ऑक्साइड को किसी हाइड्रोकार्बन से कैसे अलग किया जाता है?
  \item ट्राइफ़ेनिलफ़ॉस्फ़ीन ऑक्साइड का बनना प्रेरक बल क्यों है?
  \item HWE अभिक्रिया इस लक्ष्य के लिए उपयुक्त क्यों नहीं होती?
\end{enumerate}

\medskip
\noindent\textbf{भाग IV --- ऐल्काइन मार्ग।}
\begin{enumerate}[resume]
  \item कौन-सी ऐल्काइन, हाइड्रोजनित होकर, लक्ष्य देती है, और किस उत्प्रेरक से?
  \item डेक-1-आइन और एक हैलोऐल्केन से वह ऐल्काइन बनाइए: कौन-सा हैलाइड, कौन-सा क्षारक?
  \item हाइड्रोजनन ऐल्कीन पर क्यों रुकना चाहिए, और उत्प्रेरक इसे कैसे सुनिश्चित करता है?
  \item दोनों मार्गों की तुलना कीजिए: चरणों की संख्या, ज्यामिति नियंत्रण, उपोत्पाद।
  \item हाइड्रोजनन चरण की परमाणु मितव्ययिता क्या है?
  \item विटिग मार्ग से फ़ेरोमोन के प्रति ग्राम सह-उत्पादित ट्राइफ़ेनिलफ़ॉस्फ़ीन ऑक्साइड का द्रव्यमान
        बताइए।
\end{enumerate}
\end{problem}
